%% parts/p07/calculi-and-derivability-supplement.tex
% Necessity: the exact inherited calculus chapter replaces the overlapping
% historical exposition but does not contain the separately attributed
% registry comparison, translation results, or apparatus tables.
% Definition: this separately authored addition retains the historical
% lines 79--90 soundness statement and revises the former lines 132--318
% reader material.  The revision updates active-sheet routes, supplies the
% complete LK--ND/completeness proofs, and adds the A-117 correspondence figure.
% Standing: it is an administrative reader supplement, not inherited source,
% an admitted theorem family, or evidence that any register row is complete.
% Nonredundancy: the source supplies its own LK/LJ rules and cut results; this
% addition instead distinguishes registered presentations and their routes.
% Use: the Part VII driver routes it after all eight exact opening-chapter
% wrappers; it is not represented as restored source.

\section{Registered-calculus additions}

% BEGIN HISTORICAL SLICE 79-90
\begin{hbtheorem}{thm:p07-lkd-soundness}{sketched=annex:p7.calculi.1}{Soundness of the defined-negation sequent calculus}
Let \(\Gamma,\Delta\) be finite contexts.  If
\(\Gamma\derives_{\mathsf{LK}_{\mathrm d}}\Delta\), then every
classical structure and assignment satisfying every formula occurrence in
\(\Gamma\) satisfies at least one formula occurrence in \(\Delta\).  For an
empty succedent, the conclusion says that no structure and assignment
satisfies every antecedent formula.
\end{hbtheorem}

Semantic soundness of \(\mathsf{LJ}\) is not inferred from this classical
result.  Its Kripke semantics and adequacy belong to the Part-8
intuitionistic-logic entry.
% END HISTORICAL SLICE 79-90

% BEGIN HISTORICAL DERIVATIVE 132-318
The citation-only comparison rows below use Girard's IMLL presentation
\cite{Girard1987} and the basic ordered Lambek presentation surveyed in
\cite{GalatosJipsenKowalskiOno2007}; neither citation supplies a P7 registry
sheet.

\begin{table}[!htbp]
\centering
\begingroup\scriptsize
\setlength{\tabcolsep}{3pt}
\begin{tabular}{@{}p{.18\textwidth}p{.12\textwidth}p{.12\textwidth}
  p{.12\textwidth}p{.12\textwidth}p{.25\textwidth}@{}}
\toprule
Exact calculus & Exchange & Weakening & Contraction & Cut displayed? &
Admissibility evidence / scope \\
\midrule
Canonical \(\mathsf{LK}\), \texttt{sheet:lk} & Built into both multiset
contexts & Primitive left and right & Primitive left and right & Primitive &
Not established by the inherited cut-reduction sketch migrated from source
Part VIII, whose logical and identity rules are those of
\(\mathsf{LK}_{\mathrm d}\). \\
\(\mathsf{LK}_{\mathrm d}\), \texttt{sheet:lk-defneg} & Built into both
multiset contexts & Primitive left and right & Primitive left and right &
Primitive & Sketched in Chapter~\ref{ch:cut-elimination}; equality and
primitive negation are outside its scope. \\
\(\mathsf{LJ}\), \texttt{sheet:lj} & Built into the antecedent multiset &
Primitive on the left; no right weakening in the exact-single-succedent
syntax & Primitive on the left; no right contraction & Primitive & Sketched
in Chapter~\ref{ch:cut-elimination}. \\
Girard's commutative IMLL, \texttt{sheet:imll} & Built into multiset context &
Absent & Absent & Displayed in the cited presentation & Cut elimination is
citation-only here; the registry supplies the exact displayed rules. \\
Associative Lambek calculus \(\mathsf L\) with ordered antecedents & Absent &
Absent & Absent & Displayed in the selected basic sequent presentation &
Cut admissibility is citation-only; the ordered context, rather than a
multiset, blocks exchange. \\
Relevant implication \(\mathbf R\), \texttt{sheet:relevance-r} & Not a sequent
context on the selected Hilbert sheet & Not a primitive context rule & Not a
primitive context rule & No sequent Cut on the selected sheet & The row records
the exact Hilbert presentation and makes no family-wide structural inference. \\
Commutative affine ILL, \texttt{sheet:affine-ill} & Built into multiset context &
Primitive antecedent weakening & Absent & Inherited from \texttt{sheet:imll} &
Cut elimination is citation-carried for the selected affine presentation
\cite{Lincoln1995}; the sheet itself does not prove it. \\
\bottomrule
\end{tabular}
\endgroup
\caption{Structural rules by exact presentation.  Primitive
Cut and admissibility are separate columns.  Relevance logics have no
family-wide structural profile: contraction and mingle vary by system and
presentation. Their concrete Routley--Meyer/system-specific treatment remains
Part-8-owned; the selected \(\mathbf R\) row is explicitly Hilbert-style rather
than forced into a sequent profile.}
\label{tab:p07-structural-rules}
\end{table}

The relevance and affine rows now point to concrete registry presentations.
The \(\mathbf R\) row records that the selected sheet has no sequent-context
structural rules, while the affine row adds exactly weakening to IMLL and keeps
contraction absent. Cut display and cut-admissibility evidence remain separate
columns.

\section{Completeness and presentation translations}

\begin{hbmethod}{def:p07-lk-completeness-lookup}{Finite-sequent completeness statement guide}
For the propositional restriction of canonical
\hyperref[sheet:lk]{\texttt{sheet:lk}}, finite-multiset
\(\Gamma\Rightarrow\Delta\) is derivable iff every valuation satisfying all
members of \(\operatorname{supp}(\Gamma)\) satisfies some member of
\(\operatorname{supp}(\Delta)\); empty contexts are included.  For equality-free
first-order canonical \(\mathsf{LK}\), the same finite-sequent equivalence
over all structures and assignments holds assuming the Axiom of Choice.
The complete local proof of both statements is
Theorem~\ref{thm:p07-canonical-lk-completeness}: it uses the proved
LK--natural-deduction translations and the already proved natural-deduction
completeness theorems.  It uses Cut.  The cut-elimination sketch below is about
\(\mathsf{LK}_{\mathrm d}\), so it does not establish cut-free completeness
for canonical \(\mathsf{LK}\).
\end{hbmethod}

The inherited Theorems~\ref{thm:prop-lk-completeness} and
\ref{thm:fol-lk-completeness} retain their source proof-sketch statuses.  The
new theorem below supplies their omitted translation subroutine without
altering those source labels.  Propositional and first-order soundness remain
the proved Theorems~\ref{thm:prop-lk-soundness} and
\ref{thm:fol-lk-soundness}.

\begin{hbdefinition}{def:p07-defined-negation-translation}{Translation from canonical LK to the defined-negation language}
For equality-free formulas, define \(d\) recursively by
\[
 d(\bot)=\bot,\qquad d(\neg\varphi)=d(\varphi)\to\bot,
\]
and let \(d\) commute with every other propositional connective and with
\(\forall,\exists\), preserving variables and terms.  Apply \(d\)
occurrencewise to contexts.  This is a metalinguistic formula translation,
not a connective or an assertion of formula identity.  It is needed because
canonical \(\mathsf{LK}\) has primitive negation while
\(\mathsf{LK}_{\mathrm d}\) does not; the remaining clauses prevent it from
changing the common equality-free grammar.
\index{defined-negation translation}
\end{hbdefinition}

\begin{hbtheorem}{thm:p07-lk-presentation-equivalence}{sketched=annex:p7.calculi.2}{Canonical/defined-negation LK equivalence}
\par\noindent
For every finite equality-free sequent \(\Gamma\Rightarrow\Delta\),
\[
 \Gamma\derives_{\mathsf{LK}}\Delta
 \quad\text{iff}\quad
 d(\Gamma)\derives_{\mathsf{LK}_{\mathrm d}}d(\Delta).
\]
The theorem compares existence of derivations only.  It asserts no identity
of derivation trees and no preservation of height, size, Cut occurrences, or
cut-freeness.
\end{hbtheorem}

\begin{hbdefinition}{def:p07-sequent-nd-translation}{Finite-succedent translation to classical natural deduction}
For a finite multiset \(\Delta\), let \(\bigvee\Delta\) be a fixed associated
disjunction of the distinct formulas in \(\operatorname{supp}(\Delta)\), in
the registry's formula order, and set
\(\bigvee\varnothing\defeq\bot\).  On the assumption side use the finite set
\(\operatorname{supp}(\Gamma)\).  This convention is required to translate
a multi-succedent sequent into the single-conclusion judgment of
\hyperref[sheet:nd-fol]{\texttt{sheet:nd-fol}}; it forgets multiplicity and
therefore is not a translation of proof objects.
\index{sequent!translation to natural deduction}
\end{hbdefinition}

\begin{hbtheorem}{thm:p07-lk-nd-equivalence}{proved=annex:p7.calculi.3}{Derivability equivalence of LK and classical natural deduction}
For every finite equality-free sequent \(\Gamma\Rightarrow\Delta\), canonical
\(\mathsf{LK}\) derives \(\Gamma\Rightarrow\Delta\) iff the classical
natural-deduction calculus on
\hyperref[sheet:nd-fol]{\texttt{sheet:nd-fol}}, restricted to the
equality-free language, derives \(\bigvee\Delta\) from
\(\operatorname{supp}(\Gamma)\).  The propositional single-succedent case is
the cited Theorem~\ref{thm:p02-calculi-derivability-equivalence}.
\end{hbtheorem}

% Apparatus A-117.
\begin{figure}[!htbp]
\centering
\begin{minipage}[t]{.43\textwidth}
\centering\small
\textbf{Natural deduction}\par\medskip
\[
 \begin{array}{c}
 \Gamma\vdash A\to B\qquad\Delta\vdash A\\[-2pt]
 \hline
 \Gamma\cup\Delta\vdash B\quad(\to E)
 \end{array}
\]
\end{minipage}
\hfill
\begin{minipage}[t]{.53\textwidth}
\centering\small
\textbf{Canonical LK translation}\par\medskip
\[
 \begin{array}{c}
 A\Rightarrow A,B\qquad A,B\Rightarrow B\\[-2pt]
 \hline
 A,A\to B\Rightarrow B\quad(\to L)
 \end{array}
\]
\[
 \begin{array}{c}
 \Delta\Rightarrow A\qquad A,A\to B\Rightarrow B\\[-2pt]
 \hline
 \Delta,A\to B\Rightarrow B\quad(\mathrm{Cut})\\[3pt]
 \Gamma\Rightarrow A\to B\qquad
 \Delta,A\to B\Rightarrow B\\[-2pt]
 \hline
 \Gamma,\Delta\Rightarrow B\quad(\mathrm{Cut})
 \end{array}
\]
\end{minipage}
\caption{One rule-level sequent--natural-deduction correspondence.  Natural
deduction packages implication use as \(\to E\); canonical LK first exposes
the implication by \(\to L\) and then composes the supplied derivations by
Cut.  The annex proves the translation for every structural, propositional,
and quantifier rule, including empty succedents and eigenvariables.}
\label{fig:p07-sequent-nd-correspondence}
\end{figure}

\begin{hbtheorem}{thm:p07-canonical-lk-completeness}{proved=annex:p7.calculi.4}{Finite-sequent completeness of canonical LK}
For every finite propositional sequent \(\Gamma\Rightarrow\Delta\), canonical
\(\mathsf{LK}\) derives the sequent exactly when it is valid under every
valuation.  For every finite equality-free first-order sequent, the analogous
equivalence holds under every structure and assignment, assuming the Axiom of
Choice for an arbitrary signature; for a finite or countable signature the
countable-language completeness route needs no Choice.  Empty antecedents and
empty succedents are included, and no cut-free-completeness claim is made.
\end{hbtheorem}

Starting instead with \(\mathsf{LK}_{\mathrm d}\) requires the translation
\(d\).  No equality version is inferred: it would have to name both
\texttt{sheet:lk-eq} and the exact natural-deduction equality rules, and it
would not inherit the cut-elimination result below.

\begingroup\scriptsize
\begin{longtable}{@{}p{.15\textwidth}p{.17\textwidth}p{.24\textwidth}
  p{.28\textwidth}@{}}
\caption{Calculi catalog and evidence boundaries.}
\label{tab:p07-calculi-catalog}\\
\toprule
Presentation & Registry home & Derivation object & Evidence recorded here \\
\midrule
\endfirsthead
\toprule
Presentation & Registry home & Derivation object & Evidence recorded here \\
\midrule
\endhead
Canonical \(\mathsf{LK}\) & \hyperref[sheet:lk]{\texttt{sheet:lk}} & Finite
tree of two-sided finite-multiset sequents & Propositional and first-order
soundness proved in Part~2; the inherited completeness entry remains a sketch,
while Theorem~\ref{thm:p07-canonical-lk-completeness} supplies a separate local
proof, with Choice in the arbitrary-signature first-order form; no P7
cut-admissibility claim for this presentation. \\
\(\mathsf{LK}_{\mathrm d}\) &
\hyperref[sheet:lk-defneg]{\texttt{sheet:lk-defneg}} & Finite tree of
two-sided sequents with defined negation & Soundness and the translation to
canonical LK sketched here; Cut admissibility sketched below. \\
\(\mathsf{LK}_{=}\) & \hyperref[sheet:lk-eq]{\texttt{sheet:lk-eq}} &
Canonical LK tree with equality identity and replacement & Soundness proved
in Part~2; symmetric replacement uses Cut; no cut-admissibility or
completeness theorem for this equality extension is added here. \\
\(\mathsf{LJ}\) & \hyperref[sheet:lj]{\texttt{sheet:lj}} & Finite tree of
single-succedent sequents & Cut admissibility sketched below; semantic
adequacy belongs to the Part-8 IPL entry. \\
Classical natural deduction (the NK catalog alias) &
\hyperref[sheet:nd-cpl]{\texttt{sheet:nd-cpl}} and
\hyperref[sheet:nd-fol]{\texttt{sheet:nd-fol}} & Finite formula tree with
nonvacuous discharge, primitive negation, and RAA & Soundness and
completeness proved in Part~2; NK is only an alias for these exact sheets,
not a new active sheet. \\
Intuitionistic natural deduction \(\mathsf{NJ}\) &
\hyperref[sheet:nj]{\texttt{sheet:nj}} & Finite formula tree with the
sheet's exact discharge and eigenvariable conditions & Normalization sketched
in Chapter~\ref{ch:normalization-curry-howard}; no semantic adequacy
result is manufactured here. \\
Avron's GS5 hypersequent calculus &
\hyperref[sheet:gs5-avron]{\texttt{sheet:gs5-avron}} & Finite trees of
finite families of classical sequents with external structural rules & The
worked axiom-5 derivation is Definition~\ref{def:p07-gs5-derivation}; the
separate S5 adequacy and cut claim is cited in
Theorem~\ref{thm:p07-gs5-adequacy}. \\
Negri's labelled \(K\) calculus &
\hyperref[sheet:labelled-k-negri]{\texttt{sheet:labelled-k-negri}} & Finite
trees whose sequents contain world-labelled formulas and syntactic relational
atoms & Definition~\ref{def:p07-labelled-k-derivation} derives K on the sheet;
adequacy and cut admissibility remain cited in
Theorem~\ref{thm:p07-labelled-k-adequacy}. \\
Worked fragment of Wansing's display \(K\) calculus &
\hyperref[sheet:display-k-wansing]{\texttt{sheet:display-k-wansing}} &
Structured display-sequent trees with reversible structural postulates &
Definition~\ref{def:p07-display-k-derivation} prints a complete three-rule
modal-display cycle; Theorem~\ref{thm:p07-display-k-adequacy} concerns the
source's full calculus, not this fragment. \\
Stouppa's deep-inference S5 calculus &
\hyperref[sheet:deep-s5-stouppa]{\texttt{sheet:deep-s5-stouppa}} & Derivations
whose modal rules may act inside a structure context &
Definition~\ref{def:p07-deep-s5-derivation} gives the worked B-axiom instance;
S5 adequacy and cut are Theorem~\ref{thm:p07-deep-s5-adequacy}. \\
\bottomrule
\end{longtable}
\endgroup

% END HISTORICAL DERIVATIVE 132-318
